\section{Conclusion}
Before asking whether a trained forecaster behaves like population Bayes, like a library posterior or like a context-only
estimator, one should ask whether those references predict differently on the data at hand, and at what margin and
precision. For scalar autoregressive systems with a known noise scale, population knowledge about the dynamics is worth a
few hundredths of a nat and only at very short contexts; observation noise shrinks it further; with a known noise scale,
libraries of sixteen or more systems are indistinguishable from continuous priors on held-out data, while very small
libraries are not, and with an unknown noise scale mid-sized libraries are separable for some draws only. The
window opens when the noise scale is unknown, which is where we would place a neural study, after checking how the window
depends on the comparator and the prior. We offer the references, the maps and the pre-registration machinery as a
precondition for behavioural claims about in-context learners of dynamical systems.
